% === §8 Resultados y productos académicos esperados ===
% Redactado por el Redactor — Fase 6
% Contiene §8.1 Resultados esperados y §8.2 Productos académicos esperados
% (main.tex declara el \section{} e ingresa este archivo)
% Dependencias: §4.2 (OE1/OE2/OE3), §6 (fases/TRL), §7 (hitos H1–H4, meses de entrega)
% Criterios convocatoria: a(30)+b(25)+c(5)+d(20)+e(20)=100 pts, umbral ≥80

% =====================================================================
\subsection{Resultados esperados}
% =====================================================================

Los resultados del proyecto se organizan en seis familias que cubren la totalidad
del objetivo general y la cadena de valor OE1$\to$OE2$\to$OE3 (fases F1$\to$F4
de la sección~\ref{sec:metodologia}), y se liberan de forma sincronizada con los
hitos H1--H4 del plan de trabajo (sección~\ref{sec:plan_trabajo}). La
correspondencia resultados$\leftrightarrow$hitos$\leftrightarrow$objetivos es
trazable y estricta.

\medskip

\noindent\textbf{(i) Resultados conceptuales por objetivo específico.}
\begin{itemize}[leftmargin=*, itemsep=3pt]
    \item \textbf{OE1 (F1, SP1):} un instrumento de diagnóstico de aprendizaje
    profundo validado psicométricamente (IVC $\geq 0.80$, alfa de Cronbach $\geq
    0.80$); un \emph{dataset} multimodal curado de trazas de aprendizaje en
    ingeniería con gobernanza ética documentada; y un modelo computacional del
    aprendiz basado en \emph{transformers} en español con desempeño reportado
    (F1-score macro, AUC-ROC). Estos resultados constituyen la base de evidencia
    que habilita el diseño agéntico de OE2.
    \item \textbf{OE2 (F2, SP2):} una arquitectura multiagente pedagógicamente
    fundada con orquestación basada en el Model Context Protocol (MCP), un
    módulo de \emph{scaffolding} adaptativo calibrado (función
    $s(\mathbf{x})\to\{h,d,c\}$), y una tubería RAG de generación de problemas
    auténticos con evaluación formativa automatizada; integrados en un prototipo
    validado en entorno simulado (TRL~5) con métricas de desempeño agéntico
    documentadas (latencia $p_{95}$, acuerdo interevaluador $\kappa\geq 0.70$,
    tasa de alucinación factual).
    \item \textbf{OE3 (F3--F4, SP3):} evidencia cuantitativa y cualitativa del
    impacto causal del ecosistema en el aprendizaje profundo, obtenida mediante
    diseño cuasiexperimental pre--post con grupo de control (diferencia
    estadísticamente significativa $p<0.05$, tamaño del efecto $d$ de Cohen);
    un ecosistema desplegado en infraestructura híbrida con nivel de madurez
    \textbf{TRL~6} verificado y \textbf{TRL~7} proyectado mediante demostración
    operativa con un actor del sector productivo; y la documentación de
    interoperabilidad (API, MCP, LTI/Caliper) que habilita la adopción externa.
\end{itemize}

\medskip

\noindent\textbf{(ii) Generación de nuevo conocimiento.}
Se producirán \textbf{2 artículos de investigación} en revistas indexadas en
cuartiles Q1/Q2 del área de tecnología educativa e IA aplicada a la educación
—candidatos: \emph{IEEE Transactions on Learning Technologies},
\emph{Computers \& Education}, \emph{International Journal of Artificial
Intelligence in Education}—, uno centrado en el modelo computacional del
aprendiz (OE1) y otro en la evidencia causal de impacto en aula (OE3).
Adicionalmente, \textbf{1 artículo} en revista indexada (Q3/Q4 o de circulación
nacional indexada en Publindex) abordará la arquitectura agéntica con
orquestación MCP (OE2). Como soporte técnico se entregarán \textbf{2 informes
técnicos-académicos}: el de caracterización y modelo del aprendiz (H1) y el de
impacto causal cuasiexperimental (H3).

\medskip

\noindent\textbf{(iii) Transferencia tecnológica.}
Se realizarán \textbf{2 ponencias} en congresos internacionales especializados
—candidatos: \emph{International Conference on Artificial Intelligence in
Education} (AIED), \emph{Educational Data Mining} (EDM), \emph{Learning
Analytics and Knowledge} (LAK)—, un \textbf{taller (workshop)} sobre apropiación
de ecosistemas agénticos en la enseñanza de la ingeniería en un congreso
nacional o internacional, y \textbf{2 charlas de apropiación social del
conocimiento} dirigidas a comunidades académicas y al sector productivo del eje
cafetero. Estos productos materializan la diseminación del conocimiento generado
y la vinculación con actores externos.

\medskip

\noindent\textbf{(iv) Productos tecnológicos.}
El producto tecnológico central es el \textbf{ecosistema de agentes autónomos}
integrado, desplegado en infraestructura híbrida (Apple Silicon local + VPS en
nube) y validado a \textbf{TRL~6} en aula universitaria real, con proyección
documentada a \textbf{TRL~7} mediante demostración operativa con un actor del
sector productivo. Comprende: el \emph{software} del ecosistema (núcleo de
orquestación MCP, agentes, módulo de \emph{scaffolding}, generador de problemas,
evaluador formativo), el \emph{dataset} curado, el modelo del aprendiz
entrenado, y la documentación de diseño, instalación y operación. El código
fuente se publicará bajo licencia abierta permisiva en un repositorio público
(GitHub) con \emph{releases} versionados.

\medskip

\noindent\textbf{(v) Propiedad intelectual.}
Se obtendrá \textbf{1 registro de software} ante la Dirección Nacional de
Derechos de Autor de Colombia para el ecosistema agéntico integrado. Se evaluará
además la viabilidad de una \textbf{solicitud de patente} (o modelo de utilidad)
sobre el módulo de \emph{scaffolding} adaptativo computacional, dejando como
resultado el informe técnico-patrimonial de novedad y el estado del arte
patentable al cierre del proyecto.

\medskip

\noindent\textbf{(vi) Formación de recurso humano.}
El proyecto forma capital humano de investigación en todos los niveles,
maximizando el criterio de formación de la convocatoria:
\begin{itemize}[leftmargin=*, itemsep=3pt]
    \item \textbf{1 tesis de doctorado} dirigida (Doctorado en Ingeniería
    Automática y Computación, UNAL), vinculada a la validación causal del
    ecosistema (OE3); el estudiante de doctorado lidera la analítica del
    aprendizaje y el diseño cuasiexperimental en F3.
    \item \textbf{2 tesis de maestría} dirigidas (Maestría en Ingeniería
    Automática): una sobre el módulo de \emph{scaffolding} adaptativo (OE2, F2)
    y otra sobre el despliegue y transferencia tecnológica con TRL~7 (OE3, F4).
    \item \textbf{3 trabajos de grado de pregrado} dirigidos en los programas de
    Ingeniería Eléctrica/Electrónica, Ciencias de la Computación e Ingeniería de
    Software: curaduría del \emph{dataset} multimodal (F1), desarrollo del
    generador de problemas y evaluador formativo (F2), y trabajo de campo del
    estudio cuasiexperimental (F3).
    \item \textbf{Vinculación de posgrado:} participación activa del estudiante
    de doctorado y los estudiantes de maestría en la producción académica
    (coautoría de artículos y ponencias) y en el semillero de Ciencia de Datos e
    IA, fortaleciendo la cadena de formación investigativa.
\end{itemize}

\medskip

% =====================================================================
\subsection{Productos académicos esperados}
% =====================================================================

La tabla~\ref{tab:productos} estructura el portafolio de productos académicos
del proyecto, cruzando cada producto con el \textbf{criterio de evaluación} de
la Convocatoria Nacional Alianzas SIUN 2025--2027 al que aporta, el
\textbf{hito y mes de entrega} definido en el plan de trabajo
(sección~\ref{sec:plan_trabajo}), el \textbf{tipo de producto} y el
\textbf{destino} (revista, congreso o instancia) previsto. La distribución está
diseñada para cubrir la totalidad de los 100 puntos de la evaluación y superar
el umbral de aprobación ($\geq 80/100$): el criterio de \emph{calidad} (a, 30
pts) se satisface con los artículos Q1/Q2, el modelo del aprendiz y el
prototipo TRL~6; el de \emph{articulación SIUN} (b, 25 pts) con la coautoría
del GCPDS (categoría A1) y la operación en líneas del SIUN; el de
\emph{articulación externa} (c, 5 pts) con la demostración operativa ante el
actor del sector productivo; el de \emph{formación} (d, 20 pts) con las tesis y
trabajos de grado dirigidos y la vinculación de posgrado; y el de \emph{aportes
territoriales/nacionales/globales} (e, 20 pts) con los productos de apropiación
social, las charlas, el taller y el alineamiento con los ODS~4, 9 y 10.

\begin{table}[h]
\centering
\small
\caption{Portafolio de productos académicos esperados, mapeado a los criterios
de evaluación de la Convocatoria Alianzas SIUN 2025--2027 y a los hitos del plan
de trabajo (\S7). Criterios: a=Calidad (30), b=Articulación SIUN (25),
c=Articulación externa (5), d=Formación (20), e=Aportes territoriales (20).}
\label{tab:productos}
\begin{tabularx}{\textwidth}{@{}l X c c l X@{}}
\toprule
\textbf{\#} & \textbf{Producto} & \textbf{Criterio} & \textbf{Hito/Mes} & \textbf{Tipo} & \textbf{Destino} \\
\midrule

% --- Generación de conocimiento ---
1 & Artículo Q1/Q2: modelo computacional del aprendiz (OE1)
  & a, d & H4 / mes 18--20 & Artículo indexado
  & IEEE Trans.\ Learn.\ Technol.\ / Comput.\ \& Educ. \\

2 & Artículo Q1/Q2: impacto causal del ecosistema en aula (OE3)
  & a, d, e & H4 / mes 18--20 & Artículo indexado
  & Int.\ J.\ Artif.\ Intell.\ Educ.\ / Comput.\ \& Educ. \\

3 & Artículo indexado: arquitectura multiagente con orquestación MCP (OE2)
  & a, b & H4 / mes 18--20 & Artículo indexado
  & Rev.\ indexada (Publindex / Q3--Q4) \\

4 & Informe técnico-académico de caracterización y modelo del aprendiz
  & a & H1 / mes 5 & Informe técnico
  & Documento interno / SIUN \\

5 & Informe técnico de impacto causal cuasiexperimental
  & a, e & H3 / mes 16 & Informe técnico
  & Documento interno / SIUN \\

\addlinespace
% --- Productos tecnológicos ---
6 & Ecosistema de agentes autónomos integrado (prototipo TRL~6)
  & a, b & H2 / mes 12 (TRL~5) & Prototipo tecnológico
  & LabIA / aula UNAL \\

7 & Ecosistema desplegado en infraestructura híbrida (TRL~6 verificado)
  & a, b & H3 / mes 16 & Prototipo tecnológico
  & Aula real UNAL (entorno relevante) \\

8 & Repositorio de código abierto versionado + documentación de diseño/operación
  & a, b, e & H4 / mes 19 & Software / diseño
  & GitHub (licencia abierta) \\

\addlinespace
% --- Propiedad intelectual ---
9 & Registro de software del ecosistema agéntico
  & a, b & H4 / mes 19 & Registro de propiedad intelectual
  & Dir.\ Nac.\ Derechos de Autor (Colombia) \\

10 & Informe de novedad patentable del módulo de scaffolding adaptativo
  & a & H4 / mes 20 & Informe técnico-patrimonial
  & Evaluación interna / posible patente \\

\addlinespace
% --- Transferencia tecnológica / apropiación ---
11 & Ponencia internacional: ecosistema agéntico para ingeniería
  & a, c, e & H4 / mes 19--20 & Ponencia
  & AIED / EDM / LAK \\

12 & Ponencia internacional: analítica del aprendizaje y diagnóstico
  & a, e & H4 / mes 19--20 & Ponencia
  & LAK / EDM \\

13 & Taller (workshop) de apropiación de ecosistemas agénticos en ingeniería
  & b, e & H4 / mes 19 & Producto de apropiación
  & Congreso nacional / internacional \\

14 & Charla de apropiación social del conocimiento (sector productivo)
  & c, e & H4 / mes 19--20 & Producto de apropiación
  & Eje cafetero / sector software \\

15 & Acta de demostración operativa con actor del sector productivo (TRL~7)
  & a, c, e & H4 / mes 19 & Acta de demostración
  & Empresa/centro de formación técnica \\

\addlinespace
% --- Formación ---
16 & Tesis de doctorado dirigida (validación causal, OE3)
  & d & H3--H4 / mes 16--20 & Tesis de posgrado
  & Doctorado Ing.\ Automática -- UNAL \\

17 & Tesis de maestría dirigida (scaffolding adaptativo, OE2)
  & d & H2--H4 / mes 12--20 & Tesis de posgrado
  & Maestría Ing.\ Automática -- UNAL \\

18 & Tesis de maestría dirigida (despliegue y transferencia TRL~7, OE3)
  & d & H4 / mes 20 & Tesis de posgrado
  & Maestría Ing.\ Automática -- UNAL \\

19 & Trabajo de grado de pregrado (curaduría del dataset, OE1)
  & d & H1 / mes 5 & Trabajo de grado
  & Ing.\ Cs.\ Computación / Software -- UNAL \\

20 & Trabajo de grado de pregrado (generador de problemas, OE2)
  & d & H2 / mes 12 & Trabajo de grado
  & Ing.\ Software -- UNAL \\

21 & Trabajo de grado de pregrado (trabajo de campo cuasiexperimental, OE3)
  & d & H3 / mes 16 & Trabajo de grado
  & Ing.\ Eléctrica/Electrónica -- UNAL \\

22 & Vinculación de estudiantes de posgrado al semillero (coautoría)
  & b, d & Continuo & Formación investigativa
  & Semillero Cs.\ Datos e IA -- UNAL \\

\bottomrule
\end{tabularx}
\end{table}

\medskip

\noindent\textbf{Cobertura de criterios y puntaje.}
La articulación con el SIUN (criterio b, 25 pts) se materializa en la
coautoría con investigadores del GCPDS (categoría A1 de Minciencias), la
operación del proyecto bajo líneas de investigación del sistema y los productos
compartidos con el semillero. La articulación externa (criterio c, 5 pts) se
respalda con el acta de demostración operativa ante un actor del sector
productivo (producto 15) y la charla de apropiación dirigida al sector (producto
14). La formación (criterio d, 20 pts) se cubre con seis productos de formación
directa (16--21) más la vinculación continua al semillero (22), involucrando a
un estudiante de doctorado, dos de maestría y tres de pregrado. Los aportes
territoriales y globales (criterio e, 20 pts) se sustentan en el alineamiento
del ecosistema con el ODS~4 (educación de calidad), el ODS~9 (industria,
innovación e infraestructura) y el ODS~10 (reducción de desigualdades), así como
en los productos de apropiación social y la transferencia al eje cafetero
colombiano. La progresión de madurez \textbf{TRL~4 (H1)$\to$TRL~5 (H2)$\to$TRL~6
(H3)$\to$TRL~7 proyectado (H4)} queda certificada por los productos tecnológicos
6--8 y el acta de demostración 15, garantizando el cumplimiento del requisito de
transferencia tecnológica de la convocatoria y del marco institucional.
